\chapter{Background}
\label{chap:background}

Before turning to the experiments, we need to lay the groundwork and
reasoning for them. We start by defining some important concepts from
control theory, as well as introducing the system that serves as our
core example throughout (\cref{sec:preliminaries}). Then, we discuss
the current state of research around embodiment, low-dimensional
control, and in that context robustness as well (\cref{sec:motivation}).
We further reason about the relevance of this thesis, as well as comparable
papers. After that, we move on to introducing our specifically chosen
setup for our experiments, define constants used and explain
our test protocol (\cref{sec:setup}). Finally, we walk the reader
through our upcoming chapters, building on the previously introduced
information, and giving reason for certain design choices (\cref{sec:approach}).

\section{Preliminaries}
\label{sec:preliminaries}

Here, we revisit the concepts that are necessary for understanding the
following chapters. This section introduces basics of control theory, such as
linear state-space models (\cref{subsec:state-space}), the double
integrator as our example (\cref{subsec:double-integrator}), and the
linear quadratic regulator (\cref{subsec:lqr}).

\subsection{Linear State-Space Models}
\label{subsec:state-space}
 
In control theory, the evolution of an agent together with its environment is
described by a \emph{state-space model}: a rule that assigns to every point of
a state space the instantaneous rate at which the system moves through it.
Throughout this thesis we restrict attention to the class of
\emph{linear time-invariant} systems, for which this rule is linear in
both the state and the input and does not change over time~\cite[pp.~3--4]{Werner2021Control}.

Such a continuous-time system is given by
\begin{subequations}
\label{eq:ss-cont}
\begin{align}
    \dot{x}(t) &= A_c\, x(t) + B_c\, u(t), \qquad x(0) = x_0, \label{eq:ss-cont-state}\\
    y(t)       &= C\, x(t) + D\, u(t), \label{eq:ss-cont-output}
\end{align}
\end{subequations}
where $x(t) \in \R^{n}$ is the \emph{state vector}, $u(t) \in \R^{m}$ the
\emph{input vector} chosen by the agent, and $y(t) \in \R^{q}$ the
\emph{output vector}, i.e.\ the output observed and thus available to the
feedback controller. Accordingly, $A_c \in \R^{n \times n}$ is called the
\emph{system matrix}, $B_c \in \R^{n \times m}$ the \emph{input matrix},
$C \in \R^{q \times n}$ the \emph{output matrix}, and $D \in \R^{q \times m}$
the \emph{feedthrough matrix}.
 
The two equations in \eqref{eq:ss-cont} play fundamentally different roles, and
it is helpful to explain their distinction.
\Cref{eq:ss-cont-state} is a differential equation: it does not compute the
state, which is assumed to be known at every instant, but prescribes its
\emph{rate of change}. The pair $(A_c, B_c)$ therefore governs how the system
\emph{behaves}, with $A_c$ describing its autonomous dynamics and
$B_c$ describing how the input of the agent influences them. \Cref{eq:ss-cont-output}, in contrast, 
contains no derivative at all. It is an instantaneous readout, and the pair $(C, D)$ therefore
describes how the system is \emph{observed} rather than how it evolves. We set
$D = 0$ throughout, which is the standard assumption for systems in which the
actuator has no direct instantaneous influence on the sensor~\cite[p.~4]{Werner2021Control}.
 
Given an input signal $u(\cdot)$, the trajectory of \eqref{eq:ss-cont-state} is
obtained by integration and is given in closed form by the variation-of-constants
formula~\cite[p.~13, Eq.~(2.20)]{Werner2021Control}
\begin{equation}
    \label{eq:variation-of-constants}
    x(t) = e^{A_c t} x_0 + \int_{0}^{t} e^{A_c (t - s)} B_c\, u(s)\, \mathrm{d}s ,
\end{equation}
where $e^{A_c t}$ denotes the matrix exponential.
 
\paragraph{Full-state measurement.}
Until \cref{subsubsec:inflated-model}, we assume that the state is fully
measurable and therefore set $C = I_n$, which reduces
\eqref{eq:ss-cont-output} to $y(t) = x(t)$. Later,
we replace $C$ by a lift matrix $U_1$ and
inject measurement noise (\cref{subsubsec:noise-injection}),
such that the previous simplification no longer holds.
 
% -----------------------------------------------------------------------------
\subsubsection{Discretization}
\label{subsubsec:discretization}
 
Digital controllers, as in our case, do not act continuously: they sample the output at a fixed
rate and hold the computed control input constant until the next sampling point. The
learning algorithms studied in this thesis operate on such sampled trajectories
as well, so our model has to reflect this. The standard approach is the
\emph{zero-order hold} (ZOH)~\cite[p.~118]{Werner2021Control}. Fixing a sampling period $h > 0$ and writing
$x_k \coloneqq x(kh)$, the input is assumed to be piecewise constant~\cite[Eq.~(7.1)]{Werner2021Control},
\begin{equation}
    \label{eq:zoh-input}
    u(t) = u_k \qquad \text{for } t \in [kh,\, (k+1)h), \quad k \in \mathbb{N}_0 .
\end{equation}
Substituting \eqref{eq:zoh-input} into \eqref{eq:variation-of-constants} over a
single sampling interval and applying the change of variables
$\tau = (k+1)h - s$ yields the discrete-time model
\begin{equation}
    \label{eq:ss-disc}
    x_{k+1} = A_d\, x_k + B_d\, u_k ,
\end{equation}
with~\cite[p.~136, Eqs.~(7.26)--(7.28)]{Werner2021Control}
\begin{equation}
    \label{eq:zoh-matrices}
    A_d = e^{A_c h},
    \qquad
    B_d = \left( \int_{0}^{h} e^{A_c \tau} \, \mathrm{d}\tau \right) B_c .
\end{equation}

We use the implementation from SciPy~\cite{SciPy2026ContDiscrete}
to carry out this transformation in the code. Unless stated otherwise, all experiments in this
thesis are carried out in discrete time with sampling period $h = 0.01\,\mathrm{s}$, and
we drop the subscript $d$ in the following, writing simply $A$ and $B$.
 
\subsection{The Double Integrator}
\label{subsec:double-integrator}
 
To make the comparisons across all chapters
meaningful, all experiments are built on a single system: the
\emph{double integrator}, a frictionless point mass of unit mass moving along a
line under an applied force $u$~\cite[pp.~87--88]{Rao2001Naive}. Its dynamics are Newton's second law,
\begin{equation}
    \label{eq:newton}
    \ddot{p}(t) = u(t),
\end{equation}
where $p(t)$ denotes the position of the mass. We simplify by using a mass of one, so that
force and acceleration coincide numerically. Choosing position and velocity as
the state, $x = [\,p,\ \dot{p}\,]\T \in \R^{2}$, the second-order equation
\eqref{eq:newton} decomposes into the first-order system
\begin{equation}
    \label{eq:di-first-order}
    \frac{\mathrm{d}}{\mathrm{d}t} \begin{bmatrix} p(t) \\ \dot{p}(t) \end{bmatrix}
    = \begin{bmatrix} \dot{p}(t) \\ u(t) \end{bmatrix},
\end{equation}
in which the first row is the definition of velocity and the second is
\eqref{eq:newton} itself. In the notation of \eqref{eq:ss-cont} this gives
\begin{equation}
    \label{eq:di-cont-matrices}
    A_c = \begin{bmatrix} 0 & 1 \\ 0 & 0 \end{bmatrix},
    \qquad
    B_c = \begin{bmatrix} 0 \\ 1 \end{bmatrix},
    \qquad
    C = I_2, \qquad D = 0 .
\end{equation}
The entry $(A_c)_{12} = 1$ should be read as stated in \cref{subsec:state-space}:
the rate of change of position is given by the current velocity.
 
Because $A_c^{2} = 0$, the matrix exponential in
\eqref{eq:zoh-matrices} truncates after two terms, $e^{A_c h} = I + A_c h$, and
the discretization is available in closed form:
\begin{equation}
    \label{eq:di-disc-matrices}
    A_d = \begin{bmatrix} 1 & h \\ 0 & 1 \end{bmatrix},
    \qquad
    B_d = \begin{bmatrix} h^{2}/2 \\ h \end{bmatrix}.
\end{equation}
This system is the minimal nontrivial choice for our purposes. It is
second-order, so the state is not directly actuated and the controller must
account for momentum. It is \emph{controllable} for every $h > 0$, i.e.\ an input
sequence can navigate it from any initial state to any final
state~\cite[p.~127, Def.~7.2]{Werner2021Control}.%
\footnote{Controllability, with its weaker form stabilizability, and the dual notions
observability and detectability are standard structural properties of linear systems.
If the reader wishes to read up on their formal definitions, we kindly refer them
to~\cite[Chs.~3--4 and pp.~127--129]{Werner2021Control}.}
Additionally, it is not open-loop stable: an unforced discrete-time system $x_{k+1} = A x_k$ is
\emph{stable} if and only if all eigenvalues of $A$ lie strictly inside the unit
disk~\cite[p.~127, Thm.~7.1]{Werner2021Control}, i.e.\ if its \emph{spectral radius}
$\spr(A) \coloneqq \max_i \lvert \lambda_i(A) \rvert$ satisfies $\spr(A) < 1$. Here, $A_c$
has a double eigenvalue at the origin and $A_d$ correspondingly a double eigenvalue at $1$,
so $\spr(A_d) = 1$. The last property matters for the
learning experiments: an uncontrolled or badly initialized policy diverges, so
the assumption of a known stabilizing initial gain (\cref{subsec:protocol}) is a
restriction.
 
\subsection{The Linear Quadratic Regulator}
\label{subsec:lqr}

Now that we have defined our state-space model and the system we
study, it is time to introduce our controller, which is based on the
\emph{linear quadratic regulator} problem.

Here, the setting is that our system starts at an arbitrary initial state in its
state space, and the goal of our controller is to return to the zero state as
cheaply as possible. Both distance to the origin, as well as applied control
input yield certain costs, which can be combined into a single cost function, where $Q$
describes the former, while $R$
associates costs to control input. Over a finite horizon $T$, this gives the
cost~\cite[pp.~5--6, Eqs.~(1.3) and (1.5)]{Werner2021Optimal}
\begin{equation}
    \label{eq:lqr-cost-finite}
    V_T(x_0, u) = \int_{0}^{T} \left( x(t)\T Q\, x(t) + u(t)\T R\, u(t) \right) \mathrm{d}t ,
    \qquad Q = Q\T \succeq 0, \quad R = R\T \succ 0 .
\end{equation}
We notice two further restrictions here: $Q$ is required to be positive
semi-definite, meaning that deviating from the origin never reduces the cost,
and $R$ must be positive definite, such that every control input results in a cost.

As the original formula relates to a continuous problem with a fixed end time $T$,
it does not fit well with our goal of optimizing over an infinite discrete horizon.
The following formula offers exactly that~\cite[p.~137, Eq.~(7.30)]{Werner2021Control}:
\begin{equation}
    \label{eq:lqr-cost-infinite}
    J(x_0, u) = \sum_{k=0}^{\infty} \left( x_k\T Q\, x_k + u_k\T R\, u_k \right) .
\end{equation}
This sum does not necessarily converge. If the closed loop
does not decay, the costs do not vanish and $J$ tends to $+\infty$. 
Using this as our definition, $J$ becomes well defined for every gain.

\subsubsection{The Optimal Gain}
\label{subsubsec:lqr-optimal}

What makes the LQR problem special is that its solution is a static linear feedback
law~\cite[Sec.~4]{Recht2019Tour}: the optimal input at time $k$ depends only on the current state of the model.
Introducing a gain matrix $K \in \R^{m \times n}$, we write it as%
\footnote{Note that the minus in front of $K$ has been chosen as the convention for this thesis.
This turns our closed loop into the format $A - BK$, while some referenced literature
uses $A + BK$. Formulas cited have been adapted to our convention by replacing $K$ with $-K$.
}
\begin{equation}
    \label{eq:feedback-law}
    u_k = -K x_k .
\end{equation}
Our controller is thus reduced to choosing a matrix, and the
optimal controller is given by the optimal gain matrix $K^{*}$.

As we are dealing with a discrete-time system, that gain is obtained from the
\emph{discrete algebraic Riccati equation} (DARE)~\cite[Sec.~4, Eq.~(4.2)]{Recht2019Tour}
\begin{equation}
    \label{eq:dare}
    P = A\T P A - A\T P B \left( R + B\T P B \right)^{-1} B\T P A + Q ,
\end{equation}
whose stabilizing solution $P = P\T \succeq 0$ yields~\cite[Sec.~4]{Recht2019Tour}
\begin{equation}
    \label{eq:optimal-gain}
    K^{*} = \left( R + B\T P B \right)^{-1} B\T P A ,
    \qquad
    u_k^{*} = -K^{*} x_k .
\end{equation}
The matrix $P$ is additionally our
\emph{cost-to-go} matrix~\cite[Sec.~1, Eq.~(1)]{Fazel2018Global}, meaning that we can use it to calculate the optimal cost 
incurred from any state $x$ onwards with
\begin{equation}
    \label{eq:cost-to-go}
    J^{*}(x) = x\T P x .
\end{equation}

Whether \eqref{eq:dare} has such a solution at all depends on two structural
properties of the system, stabilizability and
detectability~\cite[p.~442, Thm.~8]{Kucera1972Discrete}:

\begin{theorem}[Existence and uniqueness]
\label{thm:dare-existence}
Let $(A, B)$ be stabilizable and let $(A, Q^{1/2})$ be detectable. Then the DARE
\eqref{eq:dare} has a unique symmetric positive semi-definite solution $P$
and the resulting closed loop satisfies $\spr(A - B K^{*}) < 1$.
\end{theorem}

For the double integrator, both conditions hold: $(A, B)$ is controllable and thus
stabilizable, and with $Q = I_n$ (\cref{subsec:protocol}) the pair $(A, Q^{1/2}) = (A, I_n)$
is also observable. Thus $K^{*}$ is well defined in all our experiments.

\subsubsection{The Cost as a Function of the Gain}
\label{subsubsec:lqr-cost-of-gain}

\Cref{eq:optimal-gain} solves the problem, but it solves it directly and
tells us nothing about the quality of non-optimal gains. As the learning 
algorithms of the following chapters will produce such $K$ directly, 
we need to be able to calculate the cost as a function of $K$.

Fix a gain $K$ and write the resulting closed loop as
$x_{k+1} = A_K x_k$ with $A_K \coloneqq A - BK$. Substituting
\eqref{eq:feedback-law} into \eqref{eq:lqr-cost-infinite} shows that the cost is
again a quadratic form, $J(x_0, K) = x_0\T P_K x_0$, where $P_K$ now solves
the far simpler discrete \emph{Lyapunov equation}~\cite[Sec.~2.1]{Fazel2018Global}
\begin{equation}
    \label{eq:lyapunov}
    P_K = A_K\T P_K A_K + Q + K\T R K ,
\end{equation}
which is linear in the unknown $P_K$ and has a unique solution whenever
$\spr(A_K) < 1$. This restricts us to the set of stabilizing gains
\begin{equation}
    \label{eq:stabilizing-set}
    \mathcal{K} \coloneqq \left\{ K \in \R^{m \times n} : \spr(A - BK) < 1 \right\} ,
\end{equation}
outside of which we set the cost to $+\infty$, consistent with the remark after
\eqref{eq:lqr-cost-infinite}.

One dependency remains, namely on the initial state. Evaluating the gain
based on a single $x_0$ would reward gains that happen to fit that one
point, so we instead draw the initial state from a distribution,
$x_0 \sim \mathcal{D}$, and average~\cite[Sec.~3]{Fazel2018Global}. Writing
$\Sigma_0 \coloneqq \E_{x_0 \sim \mathcal{D}}[x_0 x_0\T]$ for its second moment,
the objective~\cite[Sec.~2.1]{Fazel2018Global} becomes
\begin{equation}
    \label{eq:cost-of-gain}
    f(K) \coloneqq \E_{x_0 \sim \mathcal{D}} \left[ x_0\T P_K x_0 \right]
    = \tr \left( P_K \Sigma_0 \right) ,
\end{equation}
using the cyclic property of the trace. \Cref{eq:cost-of-gain} is the objective
that the algorithms in \cref{chap:lqr} minimize and they differ
only in how much of the model is known to them. Later, in \cref{chap:inflation,chap:perception},
the RL approaches rely on a rollout-based estimate of it to optimize over.

% ------------------------------------------

\section{Motivation}
\label{sec:motivation}

An agent acting in its environment receives observational input and responds to it with a commanded
control input.
In other words, the agent is a function in relation to its environment, and to find an optimal
behavior, we need to approximate this optimal function as closely as possible. While neural networks
are known to be universal approximators~\cite{Cybenko1989Approximation,Hornik1989Multilayer}, and
hence able to approximate any function, in the case of an embodied agent, different constraints need to
be taken into account. Ay~\cite[p.~392]{Ay2015Geometric} names additionally to external embodiment constraints
possible intrinsic constraints, and argues that such can even be ``beneficial'' for the control
structure. By the core notion of cheap design, embodiment can help produce complex
behaviors while only needing smaller controllers~\cite{Pfeifer2006HowBody}, which
inherently requires constraints. So, for such embodied agents, the goal is to find an ``embodied universal
[approximator]''~\cite[p.~3]{Montufar2015CheapControl}, which restricts the neural network
architecture to an intrinsic, unknown dimension, encoding physical properties belonging to the agent
and the environment, while disallowing impossible movements~\cite[Sec.~3]{Montufar2015CheapControl}.

Biology motivates this approach even further: here it is clear to see that by commanding a body,
some features of the control architecture can be simplified and delegated to the
embodiment~\cite{Mueller2017Morphological}. For example, McGeer's passive walker~\cite{McGeer1990Passive} clearly
shows that some movements can be produced by the body without any control input, just by choosing
the right embodiment. While obviously here being simplified to not include any control at all, it
can quite simply be extended with such~\cite[Sec.~10]{McGeer1988Stability}.
M{\"u}ller and Hoffmann~\cite[p.~4]{Mueller2017Morphological} mention such a combination, a robot
climbing up walls, whose design was inspired by geckos, indeed using energy for control, however
very little compared to an approach not exploiting this specific embodiment design. Further, it can
be seen that the complex
behaviors are not only less energy-intensive to produce~\cite[p.~4]{Mueller2017Morphological}, but also utilize muscle
synergies~\cite{DAvella2003Combinations} and thus fewer control parameters. And while it is
difficult to clearly attribute the effect of the embodiment~\cite{Zahedi2013Quantifying}, for us,
this gives reason to believe that the simplicity, and thus efficiency, comes from a
low-dimensionality of the controller~\cite[Sec.~5]{Montufar2015CheapControl}.

That this is indeed true for certain examples has been shown both in theory and practice
before~\cite{Datar2025Universal,Montufar2015Geometry,Li2018Measuring}. Li et
al.~\cite{Li2018Measuring} call this lower dimension ``intrinsic dimension'', 
which is the smallest dimension of a parameter subspace in which the model
still achieves comparable metrics. They show, for example, that
for achieving $90\,\%$ of the performance of a fully connected network on MNIST,
the dimensionality can be reduced by $99.6\,\%$. For a convolutional neural
network, it is still a $99.3\,\%$ reduction.
Others show that the ``least complex model which is able to learn
an optimal policy does learn faster''~\cite[Sec.~6]{Montufar2015Geometry}. For a neural network, being
more parameter-efficient can be done in different ways, some of which we look at in
\cref{chap:inflation}. In comparison, for LQR, where the controller is a gain matrix, the
corresponding approach is reducing its rank, for example via a nuclear-norm regularization~\cite{Recht2010Guaranteed}.
However, while these approaches are more parameter-efficient at achieving similar or slightly worse performances, it
was also observed that they tend to perform even worse when encountering
noise~\cite{Datar2025Universal}.

Robustness to noise, however, is one of the important characteristics of controllers acting under
real-world conditions~\cite{Rao2001Naive,Roy2021MachineLearning}. One approach to achieving this
in nature is the principle of ``redundancy and diversity''~\cite[p.~828]{Kitano2004Biological}, which stands
in direct contrast to our goal of reducing redundant parameterizations. Another reason for producing
robust controllers is the difficult transfer from a digital training environment to real-world
embodiments, commonly referred to as the ``reality gap''~\cite{Tobin2017Domain,Peng2018SimToReal}.
Approaches to training robust controllers typically introduce noise already during training, be it by modelling
measurement and disturbance noise explicitly, as the linear quadratic Gaussian (LQG) controller
does~\cite[p.~89]{Rao2001Naive}, or by including variations to the
dynamics~\cite{Peng2018SimToReal}.
These approaches, however, do not look at low-dimensionality, which makes sense as
Kitano~\cite[p.~833]{Kitano2004Biological} proposes that robustness cannot be increased
without sacrificing performance or accepting ``increased resource
demands''.

Before taking a look at the more complex case of RL, we start simpler by
using the simpler and better-understood LQR. This approach has been taken
before~\cite{Recht2019Tour} and does make sense given the guarantees regarding convergence and
effectiveness of policy optimization for it~\cite{Fazel2018Global,Takakura2023Structured}, as well
as the understanding of its geometry~\cite{Talebi2024Policy}.
It also offers methods which produce certain structures in its controller, like sparsity or our desired
low-dimensionality, as given by structured policy
iteration (S-PI)~\cite{Park2020Structured}. Approaches can generally be categorized into two groups if the underlying dynamics
are unknown and only data is available: estimating the model and then applying a model-based solver (indirect), or using the data
to estimate the gain matrix (direct)~\cite{Zhao2024DataEnabled}. Both the indirect, as in system identification
with uncertainty bounds~\cite{Dean2020Sample}, and the direct setting, with an $\Hinf$ constraint
that bounds the closed-loop gain from disturbance to output~\cite{Zhang2021Policy},
have been studied with regard to robustness, but, to our knowledge, not in combination with low-rank policies.

Research evaluating RL on control tasks has also shown that there is room for simplification, for
example using linear policies instead of complex, fully connected neural
networks~\cite[Sec.~4.2]{Mania2018Simple}, \cite{Rajeswaran2017Towards}. This, together with the
fact that parameters in neural networks are highly dependent, and thus redundant~\cite{Denil2013Predicting}, allows us to assume that
low-dimensionality can improve efficiency here. Examples of this exist, for example by factorizing a discretized low-rank
policy into a thin and wide matrix~\cite{Rozada2024Matrix} or by introducing an orthogonal projection inside the layer
structure~\cite{Todorov2026Orthogonal}. As a result, the dimensionality of the policy is reduced. A different approach uses ``information bottlenecks'',
which exhibit better generalization behavior in unseen environments~\cite[Sec.~4.2]{Igl2019Generalization},
which again could help us with robustness.

A clear case where low-dimensional dynamics underlie a high-dimensional observational space is in perception-based tasks.
Agents observing only images instead of the states of simple control tasks are faced with the challenge of extracting
these states from a high-dimensional, often nonlinear encoding~\cite[Sec.~1]{Watter2015Embed}. Many approaches, based on
state representation learning (SRL), exist, which use different variants to create a latent embedding of the image,
used as a state~\cite{Lesort2018State}. The most interesting approach, especially for us looking at embodied agents, is
the usage of so-called ``priors'', which restrict the latent space by using knowledge about the
environment~\cite[Sec.~2.2]{Lesort2018State}.
This is similar to Ay's argument that the embodiment actively defines the action space~\cite[Sec.~1]{Ay2015Geometric}.
Robotic priors were introduced by Jonschkowski and Brock~\cite{Jonschkowski2015Learning} in 2015, and then further
refined by Botteghi et al.~\cite{Botteghi2021LowDimensional} for continuous action spaces. A more specific version
restricted to position and velocity also exists here, training a model to encode an image into a low-dimensional
position state, from which the velocity is calculated by finite differences~\cite{Jonschkowski2017PVE}.
However, perception-based control is also affected by the problems of robust control,
as noise induced through perception directly affects the performance~\cite[Sec.~3.3]{Xu2021Perception}
and therefore requires approaches for mitigating them~\cite{Dean2020Robust}.

Our goal for this thesis is therefore to look a bit deeper at the combination of low-dimensionality
and robustness in an empirical way, motivated by the findings of Datar and Ay~\cite{Datar2025Universal}. LQR
serves as the perfect, well-investigated basis~\cite{Recht2019Tour}, offering low-rank approaches that have been proposed
and analyzed theoretically, but have not yet been applied in practice and under noise~~\cite[Sec.~4]{Park2020Structured}. Previously
investigated policies for LQR under noise in comparison did not look at low-dimensionality
specifically~\cite{Zhang2021Policy}.
For RL, low-dimensionality has been explored
previously~\cite{Agarwal2020FLAMBE,Stojanovic2024ModelFree}, and also
in comparison with robustness~\cite{Li2025AdaRL}, but this research does not take into account the embodiment we are interested in.
Closing this gap, our main interest is to compare robustness properties for low-dimensional policies for embodied
agents under increasingly more complex and noisier environments.

% ------------------------------------------

\section{Experimental Setup}
\label{sec:setup}

We now introduce with \cref{subsec:modified-loop} the modifications we apply to the
double integrator loop for our experiments, followed by the metrics in \cref{subsec:robustness}
we evaluate it on. The section is concluded with \cref{subsec:protocol}, where we state
our evaluation procedure.

\subsection{The Modified Loop}
\label{subsec:modified-loop}

\subsubsection{An Inflated Model}
\label{subsubsec:inflated-model}

To allow us to investigate a controller that has many redundant inputs and outputs, acting on a
system with low-dimensional dynamics, we first need to create such a setup for our previously introduced
double integrator. We do this by keeping the dynamics, so the state progression, like the loop in \cref{subsec:double-integrator},
but projecting the observation into a higher dimension, as well as projecting actuation down from such.

More specifically, we show the modification to our loop with two random orthonormal matrices $U_1$
and $U_2$. For an inflated input dimension $d_{\mathrm{in}}$ and output dimension $d_{\mathrm{out}}$,
the two lifts are constructed the following way:
\begin{equation}
    \label{eq:lifts}
    U_1 \in \R^{d_{\mathrm{in}} \times n}, \quad U_1\T U_1 = I_n,
    \qquad
    U_2 \in \R^{m \times d_{\mathrm{out}}}, \quad U_2 U_2\T = I_m,
\end{equation}
where $U_1$ is the orthonormal factor of the QR decomposition of a random matrix, whose
entries are sampled from the standard normal distribution. $U_2$ is the transpose of such a factor
sampled from a separately drawn matrix. Both are drawn once per seed, and then kept fixed throughout.
A single step of the loop then becomes
\begin{subequations}
\label{eq:inflated-loop}
\begin{align}
    y_k     &= U_1 x_k , \label{eq:inflated-obs}\\
    v_k     &= -K y_k , \qquad K \in \R^{d_{\mathrm{out}} \times d_{\mathrm{in}}} , \label{eq:inflated-ctrl}\\
    u_k     &= U_2 v_k , \label{eq:inflated-act}\\
    x_{k+1} &= A x_k + B u_k = A x_k + B U_2\, v_k . \label{eq:inflated-dyn}
\end{align}
\end{subequations}
Compared with \eqref{eq:ss-cont}, the output matrix thus becomes $C = U_1$, so that $q = d_{\mathrm{in}}$,
and the controller acts through the input matrix $B U_2 \in \R^{n \times d_{\mathrm{out}}}$.

This way, the controller only sees the inflated observation and acts in a higher dimension. As any actuation
effort is penalized, its weight $\tilde{R}$ acts on the inflated actuation and therefore has the dimension
$d_{\mathrm{out}} \times d_{\mathrm{out}}$. To
simplify things, we also have the state weight $\tilde{Q}$ in this lifted dimension, acting on the observation. Here and in the
following, a tilde marks the inflated counterpart of a quantity of the uninflated problem. As $U_1$ has
orthonormal columns, and we use $\tilde{Q} = I_{d_{\mathrm{in}}}$ throughout, this simplifies to
\begin{equation}
    \label{eq:lifted-state-cost}
    \begin{aligned}
    \left( U_1 x_k \right)\T \tilde{Q} \left( U_1 x_k \right)
    &= x_k\T U_1\T U_1\, x_k \\
    &= x_k\T I_n\, x_k \\
    &= x_k\T Q\, x_k , \qquad Q = I_n ,
    \end{aligned}
\end{equation}
with the uninflated $Q$ of \cref{subsec:lqr}, so it indeed still penalizes how far our mass is from zero. The cost in the modified loop then
becomes
\begin{equation}
    \label{eq:inflated-cost}
    \begin{aligned}
    J(x_0, K)
    &= \sum_{k=0}^{\infty} \left( y_k\T \tilde{Q}\, y_k + v_k\T \tilde{R}\, v_k \right) \\
    &= \sum_{k=0}^{\infty} \left( x_k\T U_1\T \tilde{Q}\, U_1 x_k + x_k\T U_1\T K\T \tilde{R} K U_1 x_k \right) ,
    \end{aligned}
\end{equation}
with the closed loop $x_{k+1} = (A - B U_2 K U_1)\, x_k$.

Throughout the thesis, we use the following synonyms for different $d_{\mathrm{in}}$ and
$d_{\mathrm{out}}$ dimensionalities, applied to the loop, which can be found in \cref{tab:inflation-sizes}.
\begin{table}[htbp]
    \centering
    \caption{Preset inflation sizes.}
    \label{tab:inflation-sizes}
    \begin{tabular}{@{}lrr@{}}
        \toprule
        Name & $d_{\mathrm{in}}$ & $d_{\mathrm{out}}$ \\
        \midrule
        \emph{base}   & 2   & 1   \\
        \emph{small}  & 10  & 5   \\
        \emph{medium} & 50  & 25  \\
        \emph{large}  & 200 & 100 \\
        \emph{pixel}  & 18\,432 & 256 \\
        \bottomrule
    \end{tabular}
\end{table}

\subsubsection{Noise Injection}
\label{subsubsec:noise-injection}

Now, the second thing we want to take a look at is robustness to
small changes or variations in the environment.
We thus introduce four different types of disturbances to be included in our model,
commonly known as noise. They are distinguished by the point at which they
affect the loop: what the controller \emph{observes} (measurement noise),
what its output actually \emph{does} (control noise), what the plant
\emph{is} (parametric uncertainty), and what else \emph{acts} on it
(external disturbance). Each of them can be activated separately, so that we
can clearly isolate the cause of an effect. Throughout this section
$k$ indexes the sampling instants, and $\sigma$ denotes the standard deviation,
or amplitude of the respective noise channel.

\paragraph{Measurement noise.}
Here we influence the sensors used to measure
the model state, such that capturing the same state twice does
not lead to the same observation $y_k$.

Let $n^{\mathrm{m}}_k$ denote the measurement noise at step $k$. Then
\begin{equation}
    \label{eq:measurement-noise}
    n^{\mathrm{m}}_k \sim \mathcal{N}\!\left(0, \sigma_{\mathrm{m}}^{2} I_{d_{\mathrm{in}}}\right),
    \qquad
    y_k = U_1 x_k + n^{\mathrm{m}}_k ,
    \qquad \sigma_{\mathrm{m}} \in \R_{+} ,
\end{equation}
drawn independently at every step.

\paragraph{Control noise.}
As the name says, acts on the applied control input. With it,
we try to model the difference between the intended movement of an actuator
and its actual action due to different factors. Here we deal with
\emph{multiplicative} noise, which scales with the input. In
detail this means that the applied controller output $\tilde{v}_k$ is obtained as
\begin{equation}
    \label{eq:control-noise}
    n^{\mathrm{c}}_k \sim \mathcal{N}\!\left(0, \sigma_{\mathrm{c}}^{2} I_{d_{\mathrm{out}}}\right),
    \qquad
    \tilde{v}_k = \left( \mathbf{1}_{d_{\mathrm{out}}} + n^{\mathrm{c}}_k \right) \odot v_k ,
\end{equation}
where $\odot$ denotes the elementwise product, so that $\sigma_{\mathrm{c}}$ is
the \emph{relative} error. Applying this to the input, the next step becomes
\begin{equation}
    \label{eq:control-noise-step}
    x_{k+1} = A\, x_k + B U_2\, \tilde{v}_k .
\end{equation}

\paragraph{Parametric uncertainty.}
This noise acts directly on the matrices of the model. Here, we either decide to keep
it time invariant, meaning $\Delta^{A}_k = \Delta^{A}_0$ and $\Delta^{B}_k = \Delta^{B}_0$ for all $k$, or to resample it
at every step. The former models a wrongly identified system, while the latter
simulates one that continuously changes. As with measurement noise, it is additive.
Applied, it acts on the system the
following way:
\begin{equation}
    \label{eq:parametric-noise}
    \left( \Delta_k^{A} \right)_{ij},\ \left( \Delta_k^{B} \right)_{ij}
    \overset{\text{i.i.d.}}{\sim} \mathcal{N}\!\left(0, \sigma_{\mathrm{p}}^{2}\right),
    \qquad
    x_{k+1} = \left( A + \Delta_k^{A} \right) x_k + \left( B + \Delta_k^{B} \right) U_2\, \tilde{v}_k .
\end{equation}
Worth noting is that resampling every
step leaves the model correct \emph{on average}, since
$\E[\Delta^{A}_k] = \E[\Delta^{B}_k] = 0$, and mainly adds variance to the trajectory. This stands
in contrast to the time-invariant version, where we have a discrepancy between
what we compute actions on and what actually happens.

\paragraph{External disturbance.}
This noise models external influences, like a wind pushing on our
agent, or rain
hitting it with some frequency. It only affects the position, so with the first unit
vector $e_1 \in \R^{n}$ and $\sigma_{\mathrm{d}}$ as the noise amplitude, we get
\begin{equation}
    \label{eq:dist-sin}
    n^{\mathrm{d}}_k = \sigma_{\mathrm{d}} \sin(0.01\,k)\, e_1 .
\end{equation}
As the disturbance acts on the true state $x_k$, it is applied before any
lifting of the observation~\eqref{eq:inflated-obs}. The loop can thus
be described as
\begin{equation}
    \label{eq:disturbance-step}
    x_{k+1} = A\, x_k + B U_2\, \tilde{v}_k + n^{\mathrm{d}}_k .
\end{equation}

\subsubsection{The Full Loop}
\label{subsubsec:full-noise}

While we, until now, have only regarded each noise type on its own,
there is also the possibility to include them all at once. Combining all the channels
with the inflated loop, a single step becomes
\begin{subequations}
\label{eq:noisy-loop}
\begin{align}
    y_k         &= U_1 x_k + n^{\mathrm{m}}_k , \label{eq:noisy-obs}\\
    v_k         &= -K y_k , \label{eq:noisy-ctrl}\\
    \tilde{v}_k &= \left( \mathbf{1}_{d_{\mathrm{out}}} + n^{\mathrm{c}}_k \right) \odot v_k , \label{eq:noisy-act}\\
    x_{k+1}     &= \left( A + \Delta_k^{A} \right) x_k + \left( B + \Delta_k^{B} \right) U_2\, \tilde{v}_k + n^{\mathrm{d}}_k . \label{eq:noisy-dyn}
\end{align}
\end{subequations}

% ---------------------------

\subsection{Metrics and Robustness}
\label{subsec:robustness}

In this thesis we try to answer two main questions. First, how well
does a trained controller perform generally, as well as under influences
unknown to it during training? Secondly, what influence does its size
have on its performance, and what defines sizes? Thus we aim to introduce
the metrics here to quantify these questions and our results.

\subsubsection{Cost and Suboptimality}
\label{subsubsec:suboptimality}

By \cref{subsubsec:lqr-optimal} no gain achieves a
lower cost \eqref{eq:cost-of-gain} than $K^{*}$ on the clean loop,
and every controller we train is therefore in the best case only as good as it
unless noise is involved.
This serves as our baseline, and we evaluate the cost produced by our controllers
in relation to it.

While we can easily calculate the infinite-horizon cost of a gain
via the Lyapunov equation \eqref{eq:lyapunov}, we can only apply it
to a closed-loop system. However, when adding noise to it,
this is not the case anymore. Additionally, the neural
networks we train in later chapters do not admit a gain-matrix-based
policy $K$, which would not allow cross-chapter comparison.

So, we instead evaluate them over a fixed horizon via rollout. Hence for
a given controller, a noise condition $c$, which is a subset of channels and an
intensity (\cref{subsubsec:noise-strength}), and a horizon $T$, we
simulate starting from a randomly drawn $x_0 \sim \mathcal{D}$.
The resulting loop \eqref{eq:noisy-loop} under $c$ thus accumulates
\begin{equation}
    \label{eq:empirical-return}
    J_T^{(i,s)}(K; c)
    = \sum_{k=0}^{T-1}
      \left( \left( U_1 x_k \right)\T \tilde{Q} \left( U_1 x_k \right) + \tilde{v}_k\T \tilde{R}\, \tilde{v}_k \right) ,
\end{equation}
where $s$ indexes the seed the run was trained under and $i$ the episode
within it.

It is important to note where $\tilde{Q}$ and $\tilde{R}$ act. As $\tilde{Q}$ penalizes the difference
from the origin, we apply it on the inflated state \emph{before} any noise is
applied. This way we can evaluate the actual position of the mass, as shown
in \eqref{eq:lifted-state-cost}. For $\tilde{R}$, we handle it differently. Here,
every actuation energy adds to the cost, irrespective of its origin (so whether the
controller commanded it or it is noise).

We first define the set of pairs $(s,i)$ whose episode stayed finite under both the
trained controller and the baseline,
\begin{equation}
    \label{eq:finite-pairs}
    \mathcal{F}_{S} \coloneqq \left\{ (s,i) \;:\; J^{(i,s)}_T(K; c) < \infty ,\ J^{(i,s)}_T(K^{*}; c) < \infty \right\} .
\end{equation}
Now we can state our first main metric, which
is the relative excess of the cost produced by a
trained controller over the cost of the optimal controller $K^{*}$
under the same condition. It is taken \emph{per episode}, meaning
an episode $(s,i)$ of a trained controller and episode $(s,i)$ of the baseline
are compared here. Both start from the same $x_0$ and are affected by
the same noise. Overall this yields
\begin{equation}
    \label{eq:suboptimality}
    S^{(i,s)}(K; c)
        = \frac{J^{(i,s)}_T(K; c) - J^{(i,s)}_T(K^{*}; c)}
                {J^{(i,s)}_T(K^{*}; c)} ,
    \qquad
    \hat{S}(K; c)
        = \frac{1}{\lvert \mathcal{F}_{S} \rvert}
          \sum_{(s,i) \in \mathcal{F}_{S}} S^{(i,s)}(K; c) ,
\end{equation}
where $S = 0$ means that the trained controller performed as well
as our baseline and $S = 0.1$ that it did $10\,\%$ worse.

\subsubsection{Robustness to Noise}
\label{subsubsec:robustness-metric}

Before introducing our second main metric, we need to define 
\begin{equation}
    \label{eq:empirical-cost}
    \hat{J}_T(K; c)
    = \frac{1}{\lvert \mathcal{F} \rvert}
      \sum_{(s,i) \in \mathcal{F}} J^{(i,s)}_T(K; c) ,
    \qquad
    \mathcal{F} = \left\{ (s,i) \;:\; J^{(i,s)}_T(K; c) < \infty \right\} ,
\end{equation}
which is the cost averaged over all episodes. In the following we will also
use its variation $\hat{J}^{(s)}_T(K; c)$, where we pool over all episodes
of a seed $s$.

Our second metric relates to robustness, and helps us evaluate how much a controller is influenced
by a noise condition $c$. It is the ratio of the test cost under noise $c$ to the test cost under
the clean loop, which we from now on will call \emph{relative degradation}. Formalized, this gives
\begin{equation}
    \label{eq:relative-degradation}
    \kappa^{(s)}(K; c) \coloneqq \frac{\hat{J}^{(s)}_T(K; c)}{\hat{J}^{(s)}_T(K; \varnothing)} ,
\end{equation}
where $\hat{J}^{(s)}_T$ is \eqref{eq:empirical-cost} containing only the episodes of seed $s$.
Thus, the closer $\kappa^{(s)}(K; c)$ is to $1$, the more robust a controller is. 
We can also compare the relative degradation to the performance of the optimal gain, which we calibrated in \eqref{eq:noise-calibration}
to specific ratios, by taking the ratio of the costs averaged over all seeds,
\begin{equation}
    \label{eq:relative-degradation-pooled}
    \kappa(K; c) \coloneqq \frac{\hat{J}_T(K; c)}{\hat{J}_T(K; \varnothing)} ,
\end{equation}
for which $\kappa_{\ell} = \kappa(K^{*}; c_{\ell})$ holds by \eqref{eq:noise-calibration}.
As a result, this gives us the possibility to compare whether a specific trained
controller is more or less robust than the baseline, or any other controller.

However, this metric is misleading when comparing a controller which has performed
very badly on the clean loop with one that performed well.
If we assume that both achieve the same $\kappa(K; c)$, this would equate
to a much larger absolute increase in the test cost for the former, compared
to the latter. To quantify this, we introduce the following metric,
\begin{equation}
    \label{eq:absolute-degradation}
    \frac{\Delta \hat{J}}{\hat{J}^{*}}(K; c) \coloneqq \frac{\hat{J}_T(K; c) - \hat{J}_T(K; \varnothing)}{\hat{J}_T(K^{*}; \varnothing)} ,
\end{equation}
and call it \emph{absolute degradation}, indicating the increase in cost caused by noise,
in units of the optimal clean cost. For $K^{*}$, it equals $\kappa_{\ell} - 1$, i.e.\ $0$ on the
clean loop.

\subsubsection{The Size of a Policy}
\label{subsubsec:policy-size}

We have often mentioned ``low-dimensional'' controllers until now; however,
we want to make the number of parameters comparable, not only
between different LQR controllers, but also across LQR and RL.
So it is indeed a question of size. For a neural network we can
simply count the number of weights and biases of the actor. For a gain matrix, however,
we need to find a definition that is comparable to it.

The simplest approach for a $K \in \R^{d_{\mathrm{out}} \times d_{\mathrm{in}}}$ is the number of its entries,
$d_{\mathrm{out}} \cdot d_{\mathrm{in}}$. In \cref{chap:lqr}, however, we reduce the rank of $K$.
Optimizing the storage size for $K$ of rank $r$, the singular value decomposition~\cite[p.~85, Thm.~13.1]{Werner2021Optimal}
separates it into
\begin{equation}
    \label{eq:svd}
    K = \Phi \operatorname{diag}(\chi_1, \dots, \chi_r) \Psi\T ,
    \qquad
    \Phi \in \R^{d_{\mathrm{out}} \times r}, \quad
    \Psi \in \R^{d_{\mathrm{in}} \times r} ,
\end{equation}
with $\Phi$ and $\Psi$ having orthonormal columns and $\chi_1 \geq \dots \geq
\chi_r > 0$ being the non-zero singular values. Hence, we only need to
store $d_{\mathrm{out}} \cdot r$ entries for $\Phi$, $r$ for the singular values and
$r \cdot d_{\mathrm{in}}$ for $\Psi$, which in total yields $r(d_{\mathrm{out}} + d_{\mathrm{in}} + 1)$.
For a small enough $r$ this is less, so we take the minimum of both:

\begin{definition}[Policy size]
\label{def:policy-size}
The size of a linear policy $K \in \R^{d_{\mathrm{out}} \times d_{\mathrm{in}}}$ of rank
$r = \rank(K)$ is
\begin{equation*}
    \lvert K \rvert \coloneqq \min \left\{ d_{\mathrm{out}} \cdot d_{\mathrm{in}}, \; r \, (d_{\mathrm{out}} + d_{\mathrm{in}} + 1) \right\} ,
\end{equation*}
the number of real numbers required to store it.
The size of a neural policy is its total number of parameters in the
actor, counting both fixed and trainable ones.
\end{definition}

\subsubsection{Noise Intensity}
\label{subsubsec:noise-strength}

\Cref{subsubsec:noise-injection} introduced four channels which change the system
compared to the base model. Each noise channel is separated further
into five intensities,
\begin{center}
    \emph{none} $\prec$ \emph{weak} $\prec$ \emph{standard} $\prec$
    \emph{strong} $\prec$ \emph{divergent},
\end{center}
where \emph{none} leaves the system as it is and thus serves as our
control condition. All the other intensities are calibrated by how much
they increase the cost of the optimal gain $K^{*}$ on the uninflated system.
For each channel and intensity $\ell$, we choose the standard deviation
in a way such that for the condition $c_{\ell}$ under the test protocol of \cref{subsec:protocol}
\begin{equation}
    \label{eq:noise-calibration}
    \hat{J}_T(K^{*}; c_{\ell}) = \kappa_{\ell}\, \hat{J}_T(K^{*}; \varnothing),
    \qquad
    \kappa_{\ell} \in \{1.2,\ 1.5,\ 3,\ 10\}
\end{equation}
applies for \emph{weak}, \emph{standard}, \emph{strong} and \emph{divergent}
respectively.
For $K^{*}$, the increase stays the same under inflation for all channels
except control noise (\cref{tab:noise-size-ratios}). As it acts elementwise on the inflated output, its
effect on the plant is averaged over the number of actuators: \emph{divergent} control noise raises the
cost of $K^{*}$ by a factor of $9.97$ without inflation, but only by $6.80$ and
$6.62$ at the sizes $d_{\mathrm{in}} \times d_{\mathrm{out}} = 50 \times 25$ and $200 \times 100$.
For simplicity, however, we do not vary the levels by size, but instead accept this
as a side effect.
The concrete standard deviations per channel are listed in
\cref{tab:noise-levels}.

\subsection{Test Protocol}
\label{subsec:protocol}

Having introduced our metrics, we now define a protocol for how
we train and compare the different solver approaches with these metrics.

First, we note that we always use $\tilde{Q} = I_{d_{\mathrm{in}}}$ and $\tilde{R} = I_{d_{\mathrm{out}}}$ (so $Q = I_n$ and $R = I_m$ on the uninflated system), which weighs position,
velocity and control effort equally. Initial states are drawn uniformly from the
box $[-10, 10] \times [-1, 1]$. Some solvers require an initial stabilizing gain,
which for the uninflated system is $K_0 = [\,0.1,\ 0.5\,]$, whose closed loop has $\spr(A - B K_0) \approx 0.9975$.
On an inflated system, we use
\begin{equation}
    \label{eq:initial-gain}
    \tilde{K}_0 = U_2\T K_0 U_1\T + K_{\perp} , \qquad U_2 K_{\perp} U_1 = 0 , \qquad \lVert K_{\perp} \rVert_F = 0.5 .
\end{equation}
We create the random matrix $K_{\perp}$ such that $U_2 K_{\perp} U_1 = 0$, and thus adding it does not change
the stability property of the matrix. This way, the effective gain remains $U_2 \tilde{K}_0 U_1 = K_0$,
while $\tilde{K}_0$ is full rank, which offers the possibility of rank reduction.

For each algorithm, we train ten different controllers, each under a different fixed seed,
until convergence. We define convergence for LQR solvers as follows:
\begin{definition}[Convergence]
\label{def:convergence}
Let $g_j$ be the signal that a solver observes at iteration $j$, and let $\bar{g}_j$ be its
mean over the last $N_{\mathrm{w}}$ iterations. The solver has \emph{converged} at the first
iteration $j$ for which
\begin{equation*}
    \bar{g}_j \ge \left( 1 - \varepsilon_{\mathrm{rel}} \right) \bar{g}_{j - N_{\mathrm{p}}} ,
\end{equation*}
i.e.\ once the average has not decreased by the relative tolerance $\varepsilon_{\mathrm{rel}}$
within the last $N_{\mathrm{p}}$ iterations. After an iteration with an
infinite signal (caused by $K \notin \mathcal{K}$, \eqref{eq:stabilizing-set}), the comparison
restarts, such that we never compare with a window across an infinite iteration.
\end{definition}

The signal is something the solver itself observes and optimizes for,
like the cost \eqref{eq:cost-of-gain}, an estimate, or a regularized
version of it for solvers introducing structure. Some also have
additional early-stopping criteria. They will be introduced
in detail in \cref{chap:lqr}.

For all of them, we use the best iterate by
$\bar{g}_j$ to identify the policy we want to use.

The parameters of \cref{def:convergence} for the model-based and the model-free solvers are
listed in \cref{tab:convergence-lqr}.
\begin{table}[htbp]
    \centering
    \caption{Parameters of the convergence rule (\cref{def:convergence}) for the LQR solvers.}
    \label{tab:convergence-lqr}
    \begin{tabular}{@{}lrr@{}}
        \toprule
        Parameter & Model-based solvers & Model-free solvers \\
        \midrule
        Window $N_{\mathrm{w}}$                   & 1         & 50        \\
        Patience $N_{\mathrm{p}}$                 & 50        & 100       \\
        Relative tolerance $\varepsilon_{\mathrm{rel}}$ & $10^{-4}$ & $10^{-3}$ \\
        \bottomrule
    \end{tabular}
\end{table}

The number of steps the plant progressed and solver iterations until convergence are also
counted and reported. An upper cap of $20\,000$ iterations is set to prevent looping.

For RL-based algorithms we do not use such early-stopping logic. Following our experiments in
\cref{sec:app-rl-convergence}, we observed that implementing such a rule would lead to worse
policies in the end. Instead, the best policy is achieved close to the end, which leads us
to running each algorithm for as many steps as noted in \cref{tab:rl-step-caps}, and then
taking the final policy. Every $25\,000$ environment steps throughout, we
evaluate the current policy's performance, to track its training process.
\begin{table}[htbp]
    \centering
    \caption{Environment step budget of the RL solvers per inflation size.}
    \label{tab:rl-step-caps}
    \begin{tabular}{@{}lrrrr@{}}
        \toprule
        & \emph{small} & \emph{medium} & \emph{large} & \emph{pixel} \\
        \midrule
        Budget (environment steps) & $10^{7}$ & $1.5 \cdot 10^{7}$ & $3 \cdot 10^{7}$ & $10^{7}$\textsuperscript{$*$} \\
        \bottomrule
    \end{tabular}

    \smallskip
    {\footnotesize \textsuperscript{$*$}In the camera environment of \cref{sec:pixel-env}, one step corresponds to one decision
    of the agent, which is then held for $30$ environment steps. There, the step budget is $10^{7}$ agent steps, which equals
    $3 \cdot 10^{8}$ environment steps.\par}
\end{table}

After our controllers have been trained, we evaluate them per seed for \emph{100 episodes}
each, running for \emph{$T = 2000$ steps} per episode with a sampling period of $h = 0.01\,\mathrm{s}$.
For each episode we calculate the empirical cost \eqref{eq:empirical-return} and its
suboptimality \eqref{eq:suboptimality} in comparison to the optimal gain run with the same initial state
$x_0 \sim \mathcal{D}$ and sampled noise. In the end, for each solver architecture,
we get a distribution of suboptimalities over a total of 1000 trajectories.

The inflated size is always the same between training and evaluation. For noise
experiments, however, the controller is trained on a clean system, and then exposed to one with noise
activated during testing.

On the clean system, some trained controllers might not be able to stabilize the system,
and are therefore not functional. Those are therefore removed from the metric calculations
and instead reported separately. To define a diverging controller, we differentiate two cases:
\begin{enumerate}
    \item Linear gain: the controller becomes unstable, i.e.\
    $\spr(A - B U_2 K U_1) \geq 1$.
    \item Neural network: the controller in seed $s$ is not able to control the mass,
    i.e.\ contain it in the range of initial positions at the end of any of
    its episodes, formally once $\lvert p^{(i,s)}_T \rvert > 10$ for some episode $i$.
\end{enumerate}

Under the stronger noise conditions, individual episodes might still diverge, even
though their controller was stable before. To keep our mean finite, we exclude them
from the calculation, and report them separately as the divergence rate.
\begin{equation}
    \label{eq:divergence-rate}
    D(K; c) \coloneqq \frac{1}{\lvert \mathcal{S} \rvert\, N}
    \left\lvert \left\{ (s, i) \;:\; J^{(i,s)}_T(K; c) = \infty ,\ J^{(i,s)}_T(K^{*}; c) < \infty \right\} \right\rvert ,
\end{equation}
where $\mathcal{S}$ is the set of seeds whose controller is stable on the clean loop and $N = 100$ the
number of episodes per seed.

% ------------------------------------------

\section{Approach}
\label{sec:approach}

As mentioned before, the LQR serves as the perfect starting example,
which is exactly why we start our investigation here. For our investigations, we will
use the double integrator, as mentioned in \cref{subsec:double-integrator}.

First, as a basis, we will start by comparing
the standard policy gradient (PG) approach~\cite{Fazel2018Global} with S-PI~\cite{Park2020Structured}
on the true model (\emph{TrP} and \emph{TrS}). Their relative performance to the baseline will be evaluated under the different
inflation sizes introduced in \cref{tab:inflation-sizes}. Then, the same controllers will be exposed
to the different noise channels and intensities we introduced.

Afterwards, we will take a look at the cases where we do not have access to the real model. We compare
the relative performance of both direct and indirect approaches to our baseline, which is
the optimal controller calculated via the DARE \eqref{eq:dare} on the true model. For the indirect
approaches, we will use numerical algorithms for subspace state-space system identification
(N4SID)~\cite{VanOverschee1994N4SID} to estimate the underlying model from one
trajectory. A more detailed description, as well as the exact parameters selected,
will be provided in the sections discussing the setup. With this estimated model as
the basis, we employ three different approaches, called N4SID + DARE (\emph{IdD}), N4SID + Model-based PG (\emph{IdP}) and N4SID + S-PI (\emph{IdS}).
They will be complemented by Model-free Natural Policy Gradient (\emph{MfN})~\cite{Fazel2018Global} and model-free S-PI (\emph{MfS})~\cite{Park2020Structured},
which are direct solvers and estimate the gradient from episode rollouts.
All these will also be discussed in \cref{chap:lqr}.

In \cref{chap:inflation}, we will then first introduce RL-based approaches. With proximal policy optimization
(PPO)~\cite{Schulman2017Proximal} being one of
the most used architecture variants in robotics~\cite{Rudin2022Learning,Andrychowicz2020Learning,Kaufmann2023Champion}, we choose this as our main architecture,
and apply different variations of it. First, by using only one layer and dropping the activation function and the biases,
we force the trained controller to be linear, which allows us to extract
it as a gain matrix that we can closely compare with our LQR approaches (Linear PPO, \emph{Lin}). The fully connected, standard version is named
Fully-connected PPO (\emph{MLP}), and will also be evaluated. Additionally, we will introduce structures that force the observations through a low-dimensional latent
space. Here, we will reuse the N4SID system identifier from before to identify the output matrix $\hat{C}$, before applying a change of
basis to it, to achieve orthonormality in its columns. We project
the high-dimensional observation through $\hat{C}\T$ into a lower dimension, and then train PPO with this as the input. Therefore, the name for it is State encoder + PPO (\emph{StE}).
A similar approach is used for Scalable encoder + PPO (\emph{ScE}), which is based on Jonschkowski and Brock~\cite{Jonschkowski2015Learning} and
Botteghi et al.~\cite{Botteghi2021LowDimensional}, and uses
robotic priors to derive a low-dimensional latent state via a trainable encoder. Again, the construction following this encoder is PPO. Lastly, we
implement the orthonormal bottlenecks introduced by Todorov and Sabatelli~\cite{Todorov2026Orthogonal} (Bottleneck PPO, \emph{OrB}). It consists of sampling a fixed matrix
$U_{\mathrm{b}}$ with $d_{\mathrm{b}}$ orthonormal columns, which is added at the end of a trainable feature extractor, and used to project the output into a $d_{\mathrm{b}}$-dimensional space. Like with the other
encoders, we follow it with our PPO architecture. The same feature extractor, but without the projection matrix $U_{\mathrm{b}}$,
followed by PPO, gives us FC-Extractor + PPO (\emph{FcE}).
From now on, we refer to all PPO variants by their short names.

All these low-dimensional approaches, if visualized, exhibit a ``bow-tie'' shape, which in nature can actually
improve robustness~\cite[p.~829]{Kitano2004Biological}. So this is indeed interesting to investigate. However,
first, like in \cref{chap:lqr}, we start by comparing these different approaches across the different inflation
sizes, without any noise applied, to see how they generally behave. Then, we introduce noise, and compare
our approaches under \emph{large} inflation dimensions.

Taking our results from the experiments on the artificially inflated dimensions, in \cref{chap:perception}, we create a new setup, still using the double
integrator as our underlying structure. Here, we will observe the mass via a simulated camera, and place a number of actuators
evenly distributed among angles on a half-circle. This will give us an actual high-dimensional input, as well as a high-dimensional
output, that is in reality governed by low-dimensional dynamics.

By inputting images into our existing solvers, the dimensionality increases drastically. To be able to cleanly separate effects caused
by the nonlinearity of the image construction from architecture-related effects, we will first artificially inflate our inputs and outputs 
to exactly the sizes that will later be used by the pixelated environment. This evaluation will only be done on \emph{ScE},
\emph{MLP}, \emph{OrB} and \emph{FcE}, as the other approaches either do not scale well or are not relevant to the comparison.

Then, using the same architectural constructions for these solvers, but applied to our new environment,
we compare them with the additionally introduced position-velocity encoder~\cite{Jonschkowski2017PVE} followed by PPO (\emph{PVE}),
which uses similar robotic priors to \emph{ScE}, but focuses
more on position, velocity and acceleration. It learns positional encodings from the observations, and then calculates the velocity state based
on finite differences. Here, it is quite clear that we use knowledge of the limitations of the embodiment to construct more optimal controllers,
again in a bow-tie shape. Finally, we will introduce this environment to noise as well, and observe the effects of it.